\documentclass[10pt,a4paper]{amsart}
\usepackage[margin=19mm]{geometry}
\usepackage{amsmath,amssymb,mathtools}
\usepackage[scr=rsfs,cal=euler]{mathalfa}
\usepackage{xfrac}
\usepackage[unicode,hidelinks,bookmarks=false]{hyperref}
\newcommand{\ie}{i.e.\ }
\newcommand{\cf}{cf.\ }
\newcommand{\resp}{respectively\ }
\newcommand{\Subsection}{\subsection}
\providecommand{\nobreakdash}{\nobreak\hbox{-}}
\setlength{\emergencystretch}{2em}
\multlinegap0pt
\usepackage{tikz}
\usepackage{amssymb}
\usepackage{enumerate}
\usepackage{float}
\newtheorem{proposition}{Proposition}
\newtheorem{lemma}{Lemma}
\newtheorem{theorem}{Theorem}
\title{Prime Fano $4$-folds with semi-free torus actions}
\author{Nicholas Lindsay}
\date{Source: arXiv:2510.17230v3 (CC BY 4.0)}
\begin{document}
\maketitle

\noindent\emph{Source and licence.} Prime Fano $4$-folds with semi-free torus actions. Version arXiv:2510.17230v3 (CC BY 4.0), distributed by arXiv under the Creative Commons Attribution 4.0 International licence, http://creativecommons.org/licenses/by/4.0/, as recorded in the arXiv licence metadata for this exact version and shown on the abstract page https://arxiv.org/abs/2510.17230v3. Full licence notice: \texttt{licenses/arxiv-2510.17230-CC-BY-4.0.txt}.\par\smallskip

\noindent\textit{Editorial scope.} Counted unit: the article's theorem intermediate (source0.tex lines 600--608). The article's complete original proof of it is delivered with it (source0.tex lines 609--623, one \texttt{proof} environment of 15 lines). No mathematics is written, repaired or re-derived: the statement and the proof are the article's own lines, in the article's own notation, and no retained line is substituted or re-flowed.  Retained uncounted as the source-local context the proof needs: lines 299--300; lines 338--340; lines 359--361; lines 395--397; lines 416--419; lines 562--564.  Nothing was omitted from any retained span, so the package carries no omission record.  No retained line contains a citation, so the wrapper carries no bibliography and no bibliography entry.  No retained line contains a figure environment, an image-inclusion command or drawing code, so no image is delivered and the package declares no asset dependency.  Editorial formatting only, in addition: an AMS \texttt{amsart} preamble that restates the article's own declarations and macros used by the retained lines, namely the environments \texttt{proposition}, \texttt{lemma}, \texttt{theorem} on separate counters, together with the standard packages those lines need.  The retained environments print in this document as Proposition 1, Lemma 1, Theorem 1 in order of appearance, whereas the article prints the same items with the numbering of its own full document.  The complete native source of the article as distributed in the arXiv e-print for this exact version is bundled unchanged as \texttt{sources/187397/source0.tex}.
\begin{proposition} \label{oneonecase} Let $(M,\omega)$ be a closed symplectic $8$-manifold with $b_{2}(M)=1$ having a semi-free Hamiltonian $S^1$-action such that $(d_1,d_2) = (0,0)$. Then, there is only one non-extremal fixed component $F$ symplectomorphic to $ \mathbb{CP}^1 \times \mathbb{CP}^1$, with symplectic form $4\omega_{FS} \times4\omega_{FS}$, and normal bundle $N_{F} \cong \mathcal{O}(1,1) \oplus \mathcal{O}(1,1).$
\end{proposition}

\begin{lemma} \label{zerofourcaseone}
Let $(M,\omega)$  closed symplectic $8$-manifold with $b_{2}(M)=1$ with a semi-free Hamiltonian $S^1$-action such that $(d_1,d_2)=(0,4)$ and $M^{S^1}$ consists of $M_{\max}$ and only isolated fixed points. Then if $N_{\max}$ denotes the normal bundle of $M_{\max}$, then: $M_{\max}$ is symplectomorphic to $\mathbb{CP}^2$, $$\int_{M_{\max}}c_{2}(N_{\max}) = b_{4}(F),$$ $$  \int_{M_{\max}} c_{1}^2(N_{\max})  = 1. $$ If $b_{4}(M)>1$ then $c_{1}(N_{\max})=-L$, where $L$ represents the Poincar\'{e} dual of a line. 
\end{lemma}

\begin{proposition} \label{zerofourcasetwo}
Let $M$ be a closed symplectic $8$-manifold with $b_{2}(M)=1$ with a semi-Hamiltonian $S^1$-action  such that $(d_1,d_2)=(0,4)$ and such that $M^{S^1} \setminus M_{\max}$ has a component of positive dimension. Then there exists a fixed component $\Sigma$ of dimension $2$ with $\lambda_{F}=1$. Moreover, and, if $N$ denotes the normal bundle of $M_{\max}$ $\int_{M_{\max}} c_{1}(N_{\max})^2 \in \{0,1\}$, moreover $M_{\max}$ is symplectomorphic to $\mathbb{CP}^2$. Moreover, $$\int_{M_{\max}}c_{2}(N_{\max}) = b_{4}(F)$$ and  $$\int_{\Sigma} c_{1}(M)  = 3(3- \int_{M_{\max}} c_{1}(N_{\max})^2).$$If additionally  $b_{4}(M) > 2$, then $c_{1}(N_{\max})=-L$ where $L$ is the Poincar\'{e} dual of a line in particular $\int_{M_{\max}} c_{1}(N_{\max})^2 =1$ in this case.
\end{proposition}

\begin{lemma} \label{zerosixcase}
Let $(M,\omega)$ be a closed symplectic $8$-manifold having a semi-free Hamiltonian $S^1$-action and satisfying $b_{2}(M)=1$. Assume furthermore, $(d_1.d_{2}) = (0,6)$. Then, the fixed point set has two components, $M_{\min}$ an isolated fixed point and $M_{\max}$ is symplectomorphic to $\mathbb{CP}^3$, $c_{1}(N_{\max}) = h,$ where $h \in H^{2}(\mathbb{CP}^3,\mathbb{Z})$ is the hyperplane class.
\end{lemma}

\begin{proposition} \label{twofourcase}
Let $(M,\omega)$ be a closed simply connected symplectic $8$-manifold having $b_{2}(M)=1$ and having a semi-free Hamiltonian $S^1$-action. Suppose that $(d_1,d_2) = (2,4)$. Let $N_{\min},N_{\max}$ denote the normal bundle of $M_{\min},M_{\max}$ respectively.  Then, $b_{4}(M)=1$ and $M^{S^1}=M_{\min} \cup M_{\max}$. Furthermore, $M_{\max}$ is symplectomorphic to $\mathbb{CP}^2$, $\int_{M_{\max}} c_{1}(N_{\max})^2 = 4$, $M_{\min}$ is symplectomorphic to $\mathbb{CP}^1$ and $\int_{M_{\min}} c_{1}(N_{\min}) = 3.$ 

\end{proposition}

\begin{proposition} \label{fourfourcase}
Let $(M,\omega)$ be a closed simply connected symplectic $8$-manifold having $b_{2}(M)=1$ and having a semi-free Hamiltonian $S^1$-action. Suppose that $(d_1,d_2) = (4,4)$. Let $N_{\min},N_{\max}$ denote the normal bundle of $M_{\min},M_{\max}$ respectively. Then, $M_{\min}$ and $M_{\max}$ are symplectomorphic to $\mathbb{CP}^2$,  $$\int_{M_{\min}} c_{2}(N_{\min}) + \int_{M_{\max}}c_{2}(N_{\max}) =b_{4}(M),$$ $$\int_{M_{\min}} c_{1}(N_{\min})^2 = \int_{M_{\max}} c_{1}(N_{\max})^2 =1.$$ Finally, if $b_{4}(M) > 2$ then $c_{1}(N_{\min}),c_{1}(M_{\max})$ are equal to the class$-h \in H^{2}(\mathbb{CP}^2,\mathbb{Z})$, where $h$ denotes the hyperplane class.
\end{proposition}

\begin{theorem} \label{intermediate}

 
Let $(M,\omega)$ be a closed simply connected symplectic $8$-manifold with $b_{2}(M)=1$ and a semi-free Hamiltonian $S^1$-action. Then the following hold:
\begin{itemize}
\item[a).] $M$  has Todd genus $1$.
\item[b).] All non-isolated fixed  components are symplectomorphic to either: $$ \mathbb{CP}^1,\;\mathbb{CP}^2, \; \mathbb{CP}^1 \times \mathbb{CP}^1,\; \mathbb{CP}^3,$$ with the (product of) Fubini-Study symplectic forms.
\end{itemize}
\end{theorem}

\begin{proof}

If $(d_1,d_2) = (0,0),$ By Proposition \ref{oneonecase}, in addition to the isolated fixed points $M_{\min},M_{\max}$, there is one non-extremal fixed component symplectomorphic to $N=\mathbb{CP}^1 \times \mathbb{CP}^1$, statement b). follows. Since $M_{\min}$ is isolated, the fact that $M$ has Todd genus $1$ follows from the localization formula for the Hirzebruch polynomial. 

If $(d_1,d_2) = (0,4)$ Then, since $M_{\min}$ is isolated point the Todd genus of $M$ is $1$ by the localization formula for the Hirzebruch polynomial, proving a). Statement b).  follows immediately from Lemma \ref{zerofourcaseone} and Proposition \ref{zerofourcasetwo}. 

If $(d_1,d_2) = (0,6)$ Then, since $M_{\min}$ is isolated pointthe Todd genus of $M$ is $1$ the localization formula for the Hirzebruch polynomial, proving a). By Lemma \ref{zerosixcase} $M_{\max}$ is symplectomorphic to $\mathbb{CP}^3$ and there are no non-extremal fixed components, proving b).


If $(d_1,d_2) =(2,4),$ Then, by Proposition \ref{twofourcase} $M_{\min}$ is symplectomorphic to $\mathbb{CP}^1$ a therefore the Todd genus of $M$ is $1$ by the localization formula for the Hirzebruch polynomial, proving a).  By Proposition  \ref{twofourcase}  $M_{\max}$ is symplectomorphic to $\mathbb{CP}^2$  and there are no non-extremal fixed components proving b).



If  $(d_1,d_2) =(4,4)$ Then, by Proposition \ref{fourfourcase} $M_{\min},M_{\max}$ are symplectomorphic to $\mathbb{CP}^2$  therefore  the Todd genus of $M$ is $1$ the localization formula for the Hirzebruch polynomial, proving a). Moreover by Proposition \ref{fourfourcase}, every non-extremal fixed point isolated proving b). 
\end{proof}

\end{document}
